%% expect: native native native native native native native fallback native native native native native native native native native fallback fallback fallback fallback fallback fallback native native native native native
%% reason 8: rule_thickness or axis_height in the keys, and the math fonts of this size not yet set
%% reason 18: rule_thickness or axis_height in the keys, and no arrow or baseline to read them last
%% reason 19: rule_thickness or axis_height in the keys, and no arrow or baseline to read them last
%% reason 20: rule_thickness or axis_height in the keys, and no arrow or baseline to read them last
%% reason 21: rule_thickness or axis_height in the keys, and no arrow or baseline to read them last
%% reason 22: rule_thickness or axis_height in the keys, and no arrow or baseline to read them last
%% reason 23: rule_thickness or axis_height in the keys, and no arrow or baseline to read them last
\documentclass{article}
\usepackage{amsmath}
\usepackage{tikz-cd}
\usepackage{mathptmx}
% LaTeX's \check@mathfonts sets up a size's math fonts globally when \glb@currsize is not \f@size: \glb@currsize, the math fonts and
% amsmath's \strutbox@, and the log says which fonts it found the first time, on that input line.  Real PGF's cells do it at their size,
% each on its line; at the end each arrow's rule_thickness and axis_height does it at the size around the diagram, before its labels,
% and a diagram that reads neither leaves its last cell's.
\makeatletter
\newcommand\state[1]{\typeout{[[#1: \glb@currsize, \the\ht\strutbox@+\the\dp\strutbox@]]}}
\newcommand\strutrule{\vrule width 2pt height\ht\strutbox@ depth\dp\strutbox@}
\makeatother
\newbox\unitbox
\newcommand\dims[1]{\typeout{[[#1 \the\wd\unitbox\space\the\ht\unitbox\space\the\dp\unitbox]]}}
\begin{document}
\setbox\unitbox=\hbox{\begin{tikzcd}[font=\small]
A \arrow[r, "f" {font=\scriptsize}] & B
\end{tikzcd}}\dims{1}\box\unitbox
\state{1}
% No arrow.
\setbox\unitbox=\hbox{\begin{tikzcd}[font=\footnotesize]
X
\end{tikzcd}}\dims{2}\box\unitbox
\state{2}
% A cell with a node and no text, and one whose text is on the line after its |[...]|.
\setbox\unitbox=\hbox{\begin{tikzcd}
A & |[font=\large]|
\arrow[l]
\end{tikzcd}}\dims{3}\box\unitbox
\setbox\unitbox=\hbox{\begin{tikzcd}
A \arrow[r] & |[font=\Large]|
Y
\end{tikzcd}}\dims{4}\box\unitbox
\state{4}
% No arrow, and the baseline at a cell's center or angle, which are axis_height, or at its base, which is not.
\setbox\unitbox=\hbox{\begin{tikzcd}[font=\tiny, baseline=(\tikzcdmatrixname-1-1.center)]
P \\ Q
\end{tikzcd}}\dims{5}\box\unitbox
\state{5}
\setbox\unitbox=\hbox{\begin{tikzcd}[font=\scriptsize, baseline=(\tikzcdmatrixname-2-1.30)]
P \\ Q
\end{tikzcd}}\dims{5a}\box\unitbox
\state{5a}
\setbox\unitbox=\hbox{\begin{tikzcd}[font=\small, baseline=(\tikzcdmatrixname-1-1.base)]
P \\ Q
\end{tikzcd}}\dims{6}\box\unitbox
\state{6}
% rule_thickness in the options is real PGF's formula there: the path declines it where that would set up the size.
\setbox\unitbox=\hbox{\begin{tikzcd}[row sep=10*rule_thickness]
A \arrow[d] \\ B
\end{tikzcd}}\dims{7}\box\unitbox
\setbox\unitbox=\hbox{\begin{tikzcd}[row sep=10*rule_thickness]
A \arrow[d] \\ B
\end{tikzcd}}\dims{8}\box\unitbox
\state{8}
% What amsmath's \strutbox@ does to a row that follows; \text has the diagram four times, and gather* twice that.
\begin{gather*}
  \text{\begin{tikzcd}[font=\small] Z \end{tikzcd}} \\
  c
\end{gather*}
\state{9}
% With the math fonts of the size around the diagram in place, real PGF reads rule_thickness and axis_height in that size where their
% key is used: a row's separation after each row, the last one's at \end{tikzcd} too, what \\[...] adds, a cell's inner sep after its
% text, a column's before the next cell.  Without an arrow or a baseline at a cell's center to read them after all that, the last
% formula may be one of those or a cell's, and the path declines.
Text $x$.
\setbox\unitbox=\hbox{\begin{tikzcd}[row sep=2*rule_thickness, cells={font=\small}]
A \\ B
\end{tikzcd}}\dims{10}\box\unitbox
\state{10}

$x$
\setbox\unitbox=\hbox{\begin{tikzcd}[cells={font=\small, inner sep=rule_thickness}]
A & B
\end{tikzcd}}\dims{11}\box\unitbox
\state{11}

$x$
\setbox\unitbox=\hbox{\begin{tikzcd}[cells={font=\small}]
A \\ B \\[axis_height]
\end{tikzcd}}\dims{12}\box\unitbox
\state{12}

$x$
\setbox\unitbox=\hbox{\begin{tikzcd}[row sep={axis_height}]
|[font=\small]| A \\ |[font=\small]| B
\end{tikzcd}}\dims{13}\box\unitbox
\state{13}

$x$
\setbox\unitbox=\hbox{\begin{tikzcd}[row sep=tiny, column sep=rule_thickness]
|[font=\small]| A & |[font=\small]| B
\end{tikzcd}}\dims{14}\box\unitbox
\state{14}

% \strutbox@ as the diagram leaves it, on the page.
$x$
\begin{tikzcd}[font=\small, row sep=2*rule_thickness]
Z
\end{tikzcd}\strutrule
\state{15}

% With an arrow, or a baseline at a cell's center, real PGF reads them last in the size around the diagram, and so does the path.
$x$
\setbox\unitbox=\hbox{\begin{tikzcd}[row sep=2*rule_thickness, cells={font=\small, inner sep=rule_thickness}]
A \arrow[d] \\ B \\[axis_height]
\end{tikzcd}}\dims{16}\box\unitbox
\state{16}
\setbox\unitbox=\hbox{\begin{tikzcd}[row sep=2*rule_thickness, cells={font=\small}, baseline=(\tikzcdmatrixname-2-1.center)]
A \\ B
\end{tikzcd}}\dims{17}\box\unitbox
\state{17}

% Under nodes in empty cells the node of a cell with no text, and its formula, begins at the &, its replacement or the \\ that ends
% the cell, on its line; each of these sizes is set up here for the first time.
\setbox\unitbox=\hbox{\begin{tikzcd}[nodes in empty cells, cells={font=\LARGE}]
 &
 B
\end{tikzcd}}\dims{18}\box\unitbox
\setbox\unitbox=\hbox{\begin{tikzcd}[nodes in empty cells, cells={font=\huge}]
%
 \\
 B
\end{tikzcd}}\dims{19}\box\unitbox
\setbox\unitbox=\hbox{\begin{tikzcd}[nodes in empty cells, ampersand replacement=\&, cells={font=\Huge}]
|[font=\tiny]| A \&
 \&
 B
\end{tikzcd}}\dims{20}\box\unitbox
\state{20}
\end{document}
